% 2013 湖北卷文科第 17 题：格点四边形 DEFG 与格点三角形 ABC。
\begin{tikzpicture}[font=\footnotesize, >=Stealth, line width=0.45pt, scale=0.9]
  % 网格
  \draw[gray!45, very thin] (0,0) grid (6,6);
  % 坐标轴
  \draw[->] (0,0) -- (6.55,0) node[below right] {$x$};
  \draw[->] (0,0) -- (0,6.55) node[left] {$y$};
  \node[below left] at (0,0) {$O$};
  \foreach \x in {1,...,5} {
    \draw (\x,0) -- (\x,-0.06);
    \node[below] at (\x,-0.09) {$\x$};
  }
  \foreach \y in {1,...,5} {
    \draw (0,\y) -- (-0.06,\y);
    \node[left] at (-0.09,\y) {$\y$};
  }
  % 格点三角形 ABC
  \draw[thick] (1,2) -- (1,1) -- (3,1) -- cycle;
  \fill (1,2) circle (1.6pt) (1,1) circle (1.6pt) (3,1) circle (1.6pt);
  \node[above left] at (1,2) {$B$};
  \node[below left] at (1,1) {$C$};
  \node[below] at (3,1) {$A$};
  % 格点四边形 DEFG
  \draw[thick] (2,3) -- (4,3) -- (2,5) -- (1,4) -- cycle;
  \fill (2,3) circle (1.6pt) (4,3) circle (1.6pt) (2,5) circle (1.6pt) (1,4) circle (1.6pt);
  \node[below left] at (2,3) {$D$};
  \node[right] at (4,3) {$E$};
  \node[above] at (2,5) {$F$};
  \node[left] at (1,4) {$G$};
  % 图中的其余格点
  \fill (2,4) circle (1.3pt) (3,4) circle (1.3pt) (3,3) circle (1.3pt);
\end{tikzpicture}
